\begin{figure*}[htbp]
\centering
    \includegraphics[width=\linewidth]{figs/framework.png}
    \caption{We distill aesthetic features into a feedforward Gaussian Splatting network (\textit{top}). At inference, to search for aesthetic viewpoints, we adopt a two-stage pipeline: coarse sampling to find good candidates (\textit{bottom left}) and local refinement by gradient ascent (\textit{bottom right}).}
    \label{fig:framework}
\end{figure*}

\section{Proposed Method}
\label{sec:method}

% Talk about overall framework. From sparse-view input to the aesthetic field and the two-stage viewpoint search.
\textbf{Problem Statement.} We aim to identify aesthetically pleasing viewpoints for a given 3D scene observed from sparse input views. We formulate this as finding the camera pose $\mathbf P^*$ that maximizes the predicted aesthetic score $score(\mathbf P)$ over the continuous view space:
\begin{equation}
    \mathbf P^* = \arg\max_\mathbf P score(\mathbf P).
\end{equation}
Since direct optimization is intractable as it requires back-and-forth 2D projection and evaluation, we propose to approximate this objective by learning a \textit{3D aesthetic field}, which defines a continuous mapping from camera pose to aesthetic quality while being grounded in scene geometry.

To learn this field, we \textit{distill} aesthetic knowledge from a pretrained 2D aesthetic model into a feedforward Gaussian Splatting framework~\cite{depthsplat,mvsplat} to reconstruct per-Gaussian geometry and aesthetic features from sparse input views.
The distilled field enables rendering of aesthetic features at novel viewpoints, which are evaluated by an aesthetic decoder to assess framing and composition.
To efficiently explore the 3D scene for viewpoint suggestion, we introduce a two-stage search pipeline: (1) coarse sampling of candidate viewpoints along and around the input trajectory followed by aesthetic scoring to select top candidates, and (2) gradient-based optimization for local refinement.
The overall framework is illustrated in \cref{fig:framework}.

\subsection{Distilling the 3D Aesthetic Field}\label{sec:distill}
\textbf{Feedforward 3D Gaussian Splatting.}
We build our framework upon a feedforward Gaussian Spatting model~\cite{depthsplat} that predicts per-pixel Gaussian representation from sparse input views in a single forward pass. Specifically, given input views $\{\mathbf I^i\in\mathbb{R}^{H\times W\times 3}\}_{i=1}^N$ and corresponding camera poses $\{\mathbf P^i\in\mathbb{R}^{3\times 4}\}_{i=1}^N$, a multi-view transformer extracts multi-view features $\{\mathbf F^i_{mv}\}_{i=1}^N$. The features are then fused across neighboring views using plane-sweep~\cite{collins1996space,xu2023unifying} aggregation and enriched with monocular depth cues to predict per-pixel depth, which determines the center $\boldsymbol{\mu}$ of each 3D Gaussian. A DPT head~\cite{dpt} finally regresses the remaining per-Gaussian parameters including covariance $\boldsymbol\Sigma$, opacity $\alpha$, and color $\mathbf c$. Novel view images can be rendered from the 3D Gaussians via Gaussian Splatting~\cite{kerbl20233d}.
This backbone provides a differentiable and efficient 3D representation that we extend to build the aesthetic field.

\noindent\textbf{Direct RGB Scoring and Its Instability.}
Given the feedforward Gaussian Splatting backbone, a straightforward solution for viewpoint suggestion is to render novel views and directly score them using a pretrained aesthetic model.
However, this naive approach faces two key challenges.
\textit{First}, aesthetic models are highly sensitive to small pixel variations, producing fluctuating scores even for adjacent views (see \cref{fig:nva_1}). This is because existing datasets lack annotations across nearby viewpoints in 3D, hence the models haven't seen such variations in training.
\textit{Second}, rendering artifacts at novel views can bias aesthetic predictions and misguide optimization (see \cref{fig:nva_1}).
To address them, we shift from pixel-level scoring to feature-level reasoning by distilling aesthetic representations into the 3D Gaussians.
Operating in this latent feature space enhances robustness to low-level artifacts and enforces multi-view spatial consistency, leading to smoother aesthetic variations across nearby views. See \cref{sec:nva} for further details.

\noindent\textbf{Aesthetic Feature Distillation.}
We distill high-level aesthetic representations from a pretrained aesthetic teacher model~\cite{wei2018good}.
Specifically, we select an intermediate layer as the distillation target and train our model to predict a per-Gaussian aesthetic embedding $\mathbf f_{aes}$, which can be decoded into aesthetic scores through the remaining layers of the teacher after rasterization.

As shown in \cref{fig:framework}, on top of the backbone, we introduce lightweight aesthetic modules: a CNN aesthetic encoder, an aesthetic DPT head~\cite{dpt}, and a transformer downsampler. The aesthetic encoder is from the teacher's feature extraction layers~\cite{wei2018good} and produces multi-scale aesthetic features $\{\mathbf F_{aes}^i\}_{i=1}^N$ from the input images. These features are fused with multi-view features $\{\mathbf F^i_{mv}\}_{i=1}^N$ from the backbone, and processed by the aesthetic DPT head~\cite{dpt} to regress per-Gaussian aesthetic embeddings $\mathbf f_{aes}$. Combined with other Gaussian attributes $(\boldsymbol{\mu},\boldsymbol{\Sigma},\alpha)$, these embeddings are rendered into aesthetic feature maps $\mathbf{\widehat{F}}_{pred}$ at novel views through the same rasterization pipeline as RGB rendering. To reduce storage and rasterization overhead, each $\mathbf f_{aes}$ is a compact $32$-dim vector instead of $512$-dim as the teacher feature. Since the teacher feature maps $\mathbf{F}_{gt}$ are much smaller and deeper than $\mathbf{\widehat{F}}_{pred}$, we employ a lightweight transformer downsampler to align them and obtain the final prediction $\mathbf F_{pred}$. 

Furthermore, we condition the model on camera poses at both input and novel views, since aesthetic representations are inherently viewpoint-dependent. Incorporating pose information allows the model to capture these dependencies.

\noindent\textbf{Training Objective.}
The backbone multi-view transformer, DPT head, and aesthetic encoder are kept frozen to preserve consistent geometry predictions and 2D aesthetic perception, while the additional modules are trained end-to-end. We use mean-squared error (MSE) loss on the rendered feature maps $\mathbf F_{pred}$ at hold-out views, where ground-truth $\mathbf{F}_{gt}$ are obtained by feeding ground-truth images into the pretrained aesthetic model~\cite{wei2018good}.

\subsection{Viewpoint Search Pipeline}\label{sec:search}
The learned 3D aesthetic field provides a continuous and differentiable mapping from camera pose to aesthetic score, enabling efficient search and optimization in viewpoint space.
To efficiently explore this space, we adopt a two-stage coarse-to-fine search pipeline.
In the first stage, candidate viewpoints are heuristically sampled around the input trajectories and scored to identify high-quality candidates.
In the second stage, we perform gradient ascent on the selected candidates to locally refine their poses.
The distilled aesthetic field yields a smooth score landscape (see \cref{fig:nva_1} (a)), allowing stable gradient-based optimization (see \cref{sec:grad}).
The final refined viewpoints are returned as aesthetic suggestions. The entire process is illustrated in the bottom of~\cref{fig:framework}. We now describe each stage in detail.

\noindent\textbf{Stage 1: Coarse Sampling.}
The goal of this stage is to efficiently identify promising viewpoint candidates for further refinement.
We first connect the input views into a continuous camera trajectory by interpolating both camera positions and orientations, yielding a smooth path that covers the major observed viewpoints of the scene.
Candidate viewpoints are then linearly sampled along this trajectory.
For each sample, we further generate a local neighborhood of perturbed cameras with small in-plane shifts and directional jitters, allowing local exploration around the trajectory while maintaining scene focus.
Each candidate viewpoint is rendered through the aesthetic field to obtain its aesthetic features, which are evaluated by the aesthetic decoder to produce a score.
% Thanks to the efficiency of the feedforward rasterization pipeline, all candidates can be scored efficiently.
Finally, we select the top-$K$ high-scoring candidates $\{\mathbf{P}_{cand}^k\}_{k=1}^K$ for the subsequent refinement stage. To enhance diversity, near-identical candidates are filtered out via a distance-based duplication check, ensuring that the selected set covers distinct viewpoints.

% {\color{red}Todo: (1) Talk about region-aware sampling instead of picking top-K. (2) Talk about mechanism for filtering out too identical candidates.}


\noindent\textbf{Stage 2: Gradient-based Refinement.} Starting from the candidate viewpoints obtained in Stage 1, we perform local optimization directly on camera poses to further maximize their aesthetic scores.
Specifically, each camera pose $\mathbf P_{cand}^k$ is parameterized by 3D translation and orientation, and updated through gradient ascent on the aesthetic score:
\begin{equation}
    \mathbf P_{t+1}=\mathbf P_{t} + \eta \nabla_{\mathbf P}score(\mathbf P_{t}),
\end{equation}
where $\eta$ is the step size.
In practice, we optimize a 5-dim vector comprising 3D translation and two rotation variables (yaw and pitch), as roll is rarely adjusted in typical scene captures. 
% Rotation is represented by a simple angle-based parameterization, which suffices for our local refinement setting and avoids singularity issues.
We use a fixed number of iterations for each candidate and retain the top-scoring refined viewpoints as the final suggestions.